\section{No One Is Blessed Alone}

Scripture never once shows a solitary blessed soul. Whatever the eschatological imagery is doing, it is unanimous on this: the unit of glory is plural. A multitude that no man could number, out of every nation and tribe and people and tongue, standing before the throne in white robes with palms.\endnote{Apoc 7:9, tracked verse text of the Douay--Rheims, read 2026-07-26; the vision continues to 7:17, where the Lamb leads them to the fountains of the waters of life.} A city, and a bride, and a wife of the Lamb---three corporate images for the same thing, given in a single chapter.\endnote{Apoc 21:2, 9--10, tracked verse text, read 2026-07-26. The angel says ``Come and I will shew thee the bride, the wife of the Lamb,'' and then shews him a city.} A feast, with named guests: ``many shall come from the east and the west, and shall sit down with Abraham, and Isaac and Jacob in the kingdom of heaven.''\endnote{Matt 8:11, tracked verse text, read 2026-07-26.} And, in the Epistle to the Hebrews, an address to the living that lists what they have already arrived at: ``you are come to mount Sion and to the city of the living God, the heavenly Jerusalem, and to the company of many thousands of angels, and to the church of the firstborn who are written in the heavens, and to God the judge of all, and to the spirits of the just made perfect.''\endnote{Heb 12:22--23, tracked verse text, read 2026-07-26. The rhetorical force is that the addressees have \emph{already} come to this assembly; the communion is not postponed.}

The question this section takes up is whether that plurality is part of what heaven is, or a very good thing that happens to be true of it. The question is not idle. If the company is incidental, then the whole apparatus of the communion of saints---intercession, veneration, the confidence that the dead are not lost to us---is a feature of the interim rather than of the end, and the corporate imagery is decoration. If the company is constitutive, then something has to be said about how it relates to a vision whose object is God alone.

\subsection{What the Church teaches}

The Catechism's answer is unambiguous and worth reading for its grammar. Paragraph 1024 does not say heaven includes communion; it says heaven \emph{is} communion: ``This perfect life with the Most Holy Trinity---this communion of life and love with the Trinity, with the Virgin Mary, the angels and all the blessed---is called `heaven.'\,''\endnote{\emph{Catechism of the Catholic Church} 1024, official Vatican English web text, read 2026-07-26. The Catechism's delivery renders the em-dashes as spaced hyphens; the dashes are restored here and nothing else is changed.} Paragraph 1026 repeats the identification: ``Heaven is the blessed community of all who are perfectly incorporated into Christ.''\endnote{CCC 1026, official Vatican English web text, read 2026-07-26.} And 1045, describing the consummation, makes the relation between vision and company explicit in the order of causality: ``The beatific vision, in which God opens himself in an inexhaustible way to the elect, will be the ever-flowing well-spring of happiness, peace, and mutual communion.''\endnote{CCC 1045, official Vatican English web text, read 2026-07-26. The paragraph also describes the consummation as ``the final realization of the unity of the human race, which God willed from creation.''}

The Second Vatican Council's treatment is fuller and adds the crucial claim that the communion is not interrupted by death. The pilgrims, the souls being purified, and the blessed ``in various ways and degrees are in communion in the same charity of God and neighbor''; ``the union of the wayfarers with the brethren who have gone to sleep in the peace of Christ is not in the least weakened or interrupted, but on the contrary\ldots\ is strengthened by communication of spiritual goods''; and because those in heaven ``are more closely united with Christ, they establish the whole Church more firmly in holiness.''\endnote{Second Vatican Council, \emph{Lumen gentium} n.~49, official English text at vatican.va, read 2026-07-26. The same paragraph teaches that the blessed ``do not cease to intercede with the Father for us,'' and n.~51 immediately adds the boundary: communion with those in heaven, rightly understood, ``in no way weakens, but conversely, more thoroughly enriches the latreutic worship we give to God the Father, through Christ, in the Spirit.''}

And the Catechism adds the point that makes all of this a doctrine about persons rather than about a crowd: in heaven the elect ``retain, or rather find, their true identity, their own name.''\endnote{CCC 1025, official Vatican English web text, read 2026-07-26.} Whatever union means, it does not mean dissolution into a collective. The saints are more themselves, not less.

\subsection{The counterargument, at full strength}

Against all of which stands the most awkward article in the Latin tradition on this subject, and no honest treatment can pass it by.

Aquinas asks whether the fellowship of friends is necessary for happiness. For the imperfect happiness of this life, he says, yes: a man needs friends for the sake of good action, to do them good, to delight in seeing them do good, and to be helped by them. Then:

\begin{studyframe}
``But if we speak of perfect Happiness which will be in our heavenly Fatherland, the fellowship of friends is not essential to Happiness; since man has the entire fulness of his perfection in God. But the fellowship of friends conduces to the well-being of Happiness.''\endnote{ST I-II, q.~4, a.~8, corpus, English Dominican translation, read 2026-07-26 at New Advent. He supports it from Augustine, \emph{De Genesi ad litteram} viii.25: ``the spiritual creatures receive no other interior aid to happiness than the eternity, truth, and charity of the Creator. But if they can be said to be helped from without, perhaps it is only by this that they see one another and rejoice in God, at their fellowship.'' The Augustinian locus was not independently checked for this article and is reported as Aquinas gives it.}
\end{studyframe}

And the reply to the third objection puts it as sharply as it can be put:

\begin{studyframe}
``Perfection of charity is essential to Happiness, as to the love of God, but not as to the love of our neighbor. Wherefore \emph{if there were but one soul enjoying God, it would be happy, though having no neighbor to love}. But supposing one neighbor to be there, love of him results from perfect love of God. Consequently, friendship is, as it were, concomitant with perfect Happiness.''\endnote{ST I-II, q.~4, a.~8, reply to objection 3, read 2026-07-26; emphasis added. The Latin technical term rendered ``concomitant'' is \latin{concomitans}: accompanying necessarily but not constituting.}
\end{studyframe}

One soul, alone, seeing God, would be perfectly happy. That is not a hesitation or an aside; it is the considered position of the mature \emph{Summa}, argued from the sufficiency of God, and it is exactly what a reader who thinks the communion of saints is constitutive of beatitude has to answer.

It is worth adding, in fairness to the position, that the same author does not leave the blessed indifferent to one another. In the treatise on charity he holds that the order of charity endures in heaven: ``each one will love more and reckon to be nearer to himself those who are nearer to God''---and, immediately after, that the natural bonds are not erased: ``It will however be possible in heaven for a man to love in several ways one who is connected with him, since the causes of virtuous love will not be banished from the mind of the blessed.''\endnote{ST II-II, q.~26, a.~13, corpus, read 2026-07-26 at New Advent. The article distinguishes the order of love by the good willed to another from the order by intensity of the act, and concludes that all other reasons of love are ``incomparably surpassed by that which is taken from nighness to God.'' This is, so far as the material read for this article goes, the tradition's most exact answer to the question people actually ask about seeing their families again: the bonds are not abolished, and they are re-ordered.}

\subsection{A judgment}

\begin{studybox}{Project synthesis: whether the communion of saints is constitutive of beatitude}
\textbf{Claim.} The company is constitutive of the state and not of the act, and the two claims are compatible because they answer different questions. What formally beatifies is the vision of God, and Aquinas is right that nothing else is required for the act; a single soul admitted to that vision would be blessed. But heaven, as the Church actually teaches it, is not an act in isolation---it is a state, named by the Catechism as a communion and by revelation as a city, a bride, a feast, and a numberless multitude, into which no one is admitted privately. Three considerations move this article to say that the plurality is more than an accompaniment. First, the object seen is triune: the beatitude a creature is admitted to is not the solitary self-sufficiency of a monad but the life of Father, Son, and Spirit, and it would be strange for admission to that life to be indifferent to communion. Second, the Church's own catechesis defines heaven by communion rather than describing communion as one of heaven's contents, and does so in the paragraph that follows immediately on the definition of the vision. Third, the imagery is unanimous, and unanimity in imagery is evidence about the thing imaged even when the images are not descriptions. On this reading Aquinas's ``concomitant'' is exactly right as an analysis of the beatifying act and is not a statement about how much of heaven a lone soul would have.

\smallskip
\textbf{Strongest counterargument, at full strength.} This is a distinction manufactured to save an appearance. If the fellowship of the blessed is not part of what makes them blessed, then it is an ornament, however magnificent---and Aquinas says so in terms, with an argument, from the sufficiency of God, which is the one premise a Christian cannot deny. Everything offered above is soft: the corporate imagery of the Apocalypse is ecclesial symbolism whose referent is the Church, not a report on the furniture of glory; the Catechism's ``communion of life and love'' is catechetical language, not the technical vocabulary in which ``constitutive'' has a meaning; and the appeal to the Trinity proves too much, since the same argument would make every relation in heaven constitutive of beatitude. Worse, the claim is pastorally dangerous in a specific way that the magisterium itself has flagged: \emph{Lumen gentium} 51 has to warn that communion with the saints must not weaken the worship due to God alone, and the modern sentimental version of heaven---in which the point of dying is reunion, and God is the host of the party rather than its substance---is precisely what the scholastic ``not essential'' was designed to prevent. A doctrine that makes the company constitutive has no principled defence against that drift.

\smallskip
\textbf{What would change the judgment.} An identified magisterial text asserting, in technical language, that the communion of the blessed with one another belongs to essential rather than accidental beatitude would settle it in this article's favour; nothing of the kind has been found, and CCC 1024--1026 and 1045, though strong, are catechetical rather than definitional. In the other direction, the judgment would be substantially weakened by a demonstration that the Latin tradition's ``concomitant'' has been received as a doctrinal minimum rather than as an analytic distinction---for instance by later magisterial use of ST I-II, q.~4, a.~8 in that sense. The largest gap in the evidence is that this article did not read the \emph{Summa}'s own treatment of charity in the fatherland beyond a single article, nor the Eastern tradition's ecclesiology of glory at all; either could change the balance. The judgment is held at moderate confidence.
\end{studybox}

\subsection{What follows, and what does not}

Two practical consequences follow from the teaching as it stands, and one popular consequence does not.

It follows that the dead in Christ are not out of reach. \emph{Lumen gentium} teaches the union of the living with the blessed as strengthened rather than interrupted, and the intercession of the saints as continuing; this is why a Catholic asks a saint to pray rather than merely admiring one.\endnote{\emph{Lumen gentium} nn.~49--50, read 2026-07-26; n.~50 traces the practice from ``the very first ages of the Christian religion'' and n.~51 reaffirms the decrees of Nicaea II, Florence, and Trent on the cult of the saints while urging the correction of abuses.} It also follows that the veneration of the saints is not a rival to the worship of God, and the Council says so in the same breath in which it commends the practice---a boundary that belongs beside the doctrine and not in a footnote to it.

What does not follow is the reunion scene. That the blessed know and love one another is taught; that the natural bonds are not abolished is the common teaching of the schools; that any particular person will meet any particular person, in any particular manner, at any particular age, is not taught anywhere this article has looked. The hope of seeing again those we have loved is a legitimate hope grounded in a real doctrine. The scene at the gate, with the recognitions and the explanations, is a scene, and it belongs to the register the last section of this article takes up.
